\chapter{Instruction Fetch Unit}

\section{Overview}

The Instruction Fetch Unit (IFU) is responsible for supplying a continuous stream of instructions
to the processor pipeline. Its primary responsibilities include maintaining the Program Counter
(PC), accessing the instruction cache, performing branch prediction, and delivering fetched
instructions to the decode stage.

The fetch subsystem is divided into two pipeline stages:
\begin{itemize}
  \item Fetch Stage 1 (F1)
  \item Fetch Stage 2 (F2)
\end{itemize}

This separation enables the processor to overlap instruction fetching with cache access and
branch prediction, thereby improving instruction throughput.

\begin{figure}[h]
  \centering
  \includegraphics[width=0.75\textwidth]{ch04_fetch_f1_f2.png}
  \caption{Fetch Stage 1 and Fetch Stage 2 organization}
  \label{fig:fetch-f1-f2}
\end{figure}

\section{Program Counter}

The Program Counter (PC) identifies the address of the next instruction to be fetched.
During normal execution, the PC advances sequentially by the instruction width. However, the
fetch stage may also redirect the PC based on:
\begin{itemize}
  \item Branch prediction
  \item Jump instructions
  \item Exceptions
  \item Interrupts
  \item Pipeline recovery after branch misprediction
\end{itemize}

Whenever a branch prediction is later determined to be incorrect, the recovery logic updates the
PC with the correct target address, and instruction fetching resumes from the corrected location.

\section{Fetch Stage 1 (F1)}

Fetch Stage 1 initiates instruction fetching.
Its primary responsibilities include:
\begin{itemize}
  \item Maintaining the Program Counter
  \item Accessing the Branch Target Buffer (BTB)
  \item Reading the Branch History Table (BHT)
  \item Generating prediction information
  \item Sending instruction addresses to the instruction cache
\end{itemize}

This stage attempts to predict future control flow before the instruction has even been decoded,
thereby reducing branch penalties.

\section{Instruction Cache}

The Instruction Cache stores recently accessed instructions to reduce memory access latency.
When Fetch Stage 1 supplies an instruction address, the cache checks whether the
corresponding instruction is already available.

\begin{itemize}
  \item \textbf{Cache Hit:} The instruction is returned immediately.
  \item \textbf{Cache Miss:} The memory hierarchy is accessed to retrieve the required cache line.
\end{itemize}

The instruction cache therefore plays a crucial role in sustaining a high fetch bandwidth and
minimizing pipeline stalls.

\section{Fetch Stage 2 (F2)}

Fetch Stage 2 receives the fetched instruction and associated prediction information from the
instruction cache.

\begin{figure}[h]
  \centering
  \includegraphics[width=0.45\textwidth]{ch04_fetch_unit_detail.png}
  \caption{Internal detail of the Fetch Unit: PC, +4 adder, ICache interface, and branch prediction}
  \label{fig:fetch-unit-detail}
\end{figure}

Its responsibilities include:
\begin{itemize}
  \item Buffering fetched instructions
  \item Packaging instruction metadata
  \item Forwarding branch prediction information
  \item Passing instructions to the Decode stage via the F2-to-D pipeline register
\end{itemize}

At this point, the instruction has not yet been decoded; it remains in its original binary form.
